% 2013 年北京卷文科第 8 题：正方体 ABCD-A1B1C1D1 与对角线 BD1 的三等分点 P。
% 来源 PDF：普通高考/2013/2013北京文.pdf，第 1 页（印面 789）中栏上方。
% 可见棱用实线，被遮住的 PA 类棱用虚线；BD1 为点划线，P 为其三等分点。
\begin{tikzpicture}[font=\footnotesize, line width=0.45pt]
  \coordinate (A)  at (0,0);
  \coordinate (B)  at (3,0);
  \coordinate (C)  at (4.3,1.3);
  \coordinate (D)  at (1.3,1.3);
  \coordinate (A1) at (0,3);
  \coordinate (B1) at (3,3);
  \coordinate (C1) at (4.3,4.3);
  \coordinate (D1) at (1.3,4.3);
  % 可见棱
  \draw (A) -- (B) -- (C) -- (C1) -- (B1) -- (B);
  \draw (A1) -- (B1) -- (C1) -- (D1) -- cycle;
  \draw (A1) -- (A);
  % 被遮挡的棱
  \draw[dashed] (A) -- (D) -- (C);
  \draw[dashed] (D) -- (D1);
  % 对角线 BD1 与三等分点 P
  \draw[densely dotted] (B) -- (D1);
  \coordinate (P) at ($(B)!0.3333!(D1)$);
  \fill (P) circle (1.5pt);
  \node[right] at (P) {$P$};
  % 顶点标注
  \node[below left]  at (A)  {$A$};
  \node[below right] at (B)  {$B$};
  \node[right]       at (C)  {$C$};
  \node[above left]  at (D)  {$D$};
  \node[left]        at (A1) {$A_1$};
  \node[right]       at (B1) {$B_1$};
  \node[right]       at (C1) {$C_1$};
  \node[above left]  at (D1) {$D_1$};
\end{tikzpicture}
